\section{Fundamentals}
\label{sec:fundamentals}

Throughout the paper, we use standard notation from graph theory and computational complexity, following the conventions of~\cite{Diestel2025,Cormen2022}.
All graphs under consideration are undirected and finite, unless otherwise stated.
For a graph $G = (V, E)$, we let $n = |V|$ denote the number of vertices and $m = |E|$ the number of edges.

We often deal with two distinct vertices $u, v \in V$.
The value $d(u, v)$ denotes the distance between $u$ and $v$, i.e., the length of a shortest $u$--$v$ path in $G$.
Similarly, the value $c(u, v)$ denotes the minimum size of a cut that separates $u$ and $v$.
All problems studied in this paper are defined with respect to these two distinguished vertices $u$ and $v$, which are assumed to lie in the same connected component of $G$.

We identify a path with its set of edges.
For any path $P \subseteq E$, we denote by $G \setminus P$ the graph obtained from $G$ by removing all edges of $P$.
Furthermore, for a vertex $x \in V$ and a vertex subset $S \subseteq V$, we write $\deg(x, S)$ for the number of neighbors of $x$ that belong to $S$, that is, $\deg(x, S) = |\{\,y \in S : \{x, y\} \in E\,\}|$.


The notion of a \emph{cut-path} captures a hybrid structure in a graph that combines two fundamental and seemingly opposing elements: a \emph{path} and a \emph{cut}.
Intuitively, a cut-path between two vertices $u$ and $v$ is a set of edges that simultaneously contains a path connecting $u$ to $v$ and a cut that, if removed, would disconnect $u$ from $v$.
In this sense, a cut-path ensures both \emph{communication} and \emph{control}.
Formally:

\begin{definition}[Cut-Path]
\label{def:cut-path}
Let $G = (V, E)$ be a graph, and let $u, v \in V$ be two distinct vertices.
A \emph{cut-path} between $u$ and $v$ is a set of edges $S \subseteq E$ for which there exist subsets $P, C \subseteq S$ such that:
    \begin{enumerate}
        \item $C$ is a cut separating $u$ and $v$ (i.e., removing $C$ disconnects $u$ from $v$),
        \item $P$ is a path connecting $u$ and $v$.
    \end{enumerate}
\end{definition}

\begin{definition}[Cut-Paths]
\label{def:cut-paths}
Let $G = (V, E)$ be a graph, and let $u, v \in V$ be two distinct vertices.
The set of all cut-paths between $u$ and $v$, denoted by $\CP(u, v)$, is defined as:
\[
\CP(u, v) = \{ S \mid S \text{ is a cut-path between } u \text{ and } v \}.
\]
\end{definition}

\begin{definition}[Value of Minimum Cut-Path]
\label{def:cp-value}
Let $G = (V, E)$ be a graph, and let $u, v \in V$ be two distinct vertices.
The minimum size of a cut-path between $u$ and $v$, denoted by $\cp(u, v)$, is defined as:
\[
\cp(u, v) = \min\{ |S| \mid S \in \CP(u, v) \}.
\]
\end{definition}

A cut-path $S \in \CP(u, v)$ with $|S| = \cp(u, v)$ is called a \emph{minimum cut-path} between $u$ and $v$.
The corresponding optimization problem is stated as follows.

\begin{problem}[\MinCutPath{} (Optimization Version)]
Let $G=(V,E)$ be an undirected graph and $u,v \in V$ distinct vertices; cut-paths are defined in Definition~\ref{def:cut-path}.
The optimization version of the \MinCutPath{} problem is defined as follows:

\begin{center}
\begin{tabular}{p{0.18\linewidth}p{0.75\linewidth}}
\textbf{Input:} & A graph $G=(V,E)$ and vertices $u,v \in V$. \\
\textbf{Output:} & A minimum cut-path $S^* \subseteq E$ between $u$ and $v$, together with its cardinality
\[
\cp(u,v) = |S^*|.
\]
\end{tabular}
\end{center}
\end{problem}

The following elementary bounds relate $\cp(u, v)$ to the classical quantities $c(u, v)$ and $d(u, v)$; they are used repeatedly in the subsequent sections.

\begin{lemma}[Basic Bounds]
\label{lem:basic-bounds}
Let $G = (V, E)$ be a graph, and let $u, v \in V$ be two distinct vertices. Then
\[
\max\{c(u, v), d(u, v)\} \leq \cp(u, v) \leq c(u, v) + d(u, v) - 1.
\]
In particular, if $c(u, v) = 1$ or $d(u, v) = 1$, then $\cp(u, v) = c(u, v) + d(u, v) - 1$.
\end{lemma}

\begin{proof}
Every cut-path contains a $u$--$v$ cut, which has at least $c(u, v)$ edges, and a $u$--$v$ path, which has at least $d(u, v)$ edges; this gives the lower bound.
For the upper bound, let $C$ be a minimum $u$--$v$ cut and let $P$ be a shortest $u$--$v$ path.
The set $C \cup P$ is a cut-path.
Since removing $C$ disconnects $u$ from $v$, the path $P$ contains at least one edge of $C$, and hence $|C \cup P| \leq c(u, v) + d(u, v) - 1$.
If $c(u, v) = 1$ or $d(u, v) = 1$, the lower and the upper bound coincide.
\end{proof}


Before formalizing the notion of an average $(1+\epsilon)$-approximation scheme, we briefly explain the motivation behind this concept.
In the context of \MinCutPath, exact computation may be infeasible for large random graphs.
However, many such graphs exhibit structural regularities that make near-optimal solutions achievable in polynomial time on almost all inputs.
The following definition captures this idea by requiring that the algorithm performs within a factor of $(1+\epsilon)$ of the optimum for almost every instance while still maintaining a polynomial running time.

To make the definition precise, let $I(n)$ denote the set of all valid input instances of size $n$, and let $\OPT(x)$ denote the optimal value of the objective function for a given instance $x \in I(n)$.
For minimization problems, $\OPT(x)$ is the minimum achievable value, while for maximization problems, it is the maximum achievable value.
Formally:

\begin{definition}[Average $(1+\epsilon)$-Approximation Scheme]
\label{def:approximation-scheme}
Let $\epsilon > 0$. An algorithm $A$ is called an \emph{average $(1+\epsilon)$-approximation scheme} if it meets the following criteria:
\begin{enumerate}
    \item $A$ is a polynomial-time algorithm, i.e., its running time is $O(P(n))$ for some polynomial $P(n)$.
    \item For each input size $n$, there exists a subset of inputs $I_{A}^{\mathrm{opt}}(n) \subseteq I(n)$ such that for every input $x \in I_{A}^{\mathrm{opt}}(n)$, the value $A(x)$ of the solution computed by $A$ satisfies
    \[
    \frac{A(x)}{\OPT(x)} < 1+\epsilon \quad \text{(for minimization problems)}
    \]
    \[
    \frac{\OPT(x)}{A(x)} < 1+\epsilon \quad \text{(for maximization problems)},
    \]
    where $\OPT(x)$ denotes the optimal value for the input $x$.
    \item In addition, the algorithm is average-case in the sense that
    \[
    \lim_{n \to \infty} \frac{\lvert I_{A}^{\mathrm{opt}}(n) \rvert}{\lvert I(n) \rvert} = 1.
    \]
\end{enumerate}
\end{definition}

